\scanpage{40}
\mspage{39}

Donc \deleted{\unsure{le}} on trouve

\uline{Proposition} $A$, \inserted{$k$,} $W,e,f$, \inserted{$h=ef$,} comme plus haut, $G_0$ un grp de B.T. \inserted{de hauteur $h$} sur \deleted{$A$} $k$ sur lequel $A$ opère, \deleted{\unsure{Alors $F$}} $M$ son module de Dieudonné (module libre de rang $h$ sur $W$).
$F_{G_0}^f$ est de la forme $\varpi_{G_0}$, où $\varpi\in A$, et on a $\dim G_0=\operatorname{long}A/\varpi A=v(\varpi)$.
Alors \deleted{conditions équivalentes} conditions équivalentes
\[
\begin{array}{ll}
\text{(i)}&\dim G_0=1\\
\text{(ii)}&F_{G_0}^f=\pi_{G_0},\quad\pi\ \text{une \uline{uniformisante} de }A\ \text{(uniqt déterminée !)}
\end{array}
\]
(iii) Le module de Dieudonné $M$ de $G_0$ \inserted{(sur $W$,)} avec $F_M$ et opération de $A$, est donné par un triple $(\Omega,\unsure{\nu_0|\nu},\pi)$ où $\Omega$ est un module inv. sur $A$, $0\leq\unsure{\nu_0|\nu}\leq f-1$, $\pi$ une uniformisante de $A$ par
\[
\left\{\begin{aligned}
M&=\Omega\otimes_{\mathbb Z_p}W\simeq\Omega\otimes_A(A\otimes_{\mathbb Z_p}W)\\
F_M&=\operatorname{id}_\Omega\otimes_{\mathbb Z_p}\sigma+\operatorname{pr}_\nu[(\pi-1)\otimes\sigma]\deleted{\unsure{\operatorname{pr}_\nu}}
\end{aligned}\right.
\]
\margin{(isom de $A\otimes_{\mathbb Z_p}W$-mod)}
où $\operatorname{pr}_\nu$ est la \deleted{\unsure{facteur}} projection de $M$ sur son $\nu$-ième facteur \inserted{$M_\nu$} (correspondant à la décomposition $A\otimes_{\mathbb Z_p}W\simeq A^{[0,f-1]}$).

\margin{\uline{NB} Tous ces groupes (pour $\nu$ variable) sont déduits de l'un d'entre eux (disons $\nu=0$) par \unsure{transformation} à l'aide des $f$ automorphismes de $k\simeq\mathbb F_q=\mathbb F_{p^f}$.}

\uline{NB} Revenons au cas général, où on ne suppose pas néc. $\dim G_0=1$ \inserted{(mais $G_0$ de hauteur $h\ldots$)}. Soit $G'_0$ défini comme $G_0$ par $\Omega'$, et les $r'_\nu$ $(0\leq\nu\leq f-1)$ $u'\in A^*$. On cherche les homs compatibles avec l'action de $A$ de $G_0$ dans $G'_0$, i.e. les hom de $A\otimes_{\mathbb Z_p}W$-modules de $M'$ dans $M$, compatibles avec $F_M,F_{M'}$. On trouve facilement que pour qu'il existe un tel hom qui soit $\ne0$ ($\Longleftrightarrow$ une isogénie), il f. et suffit qu'on ait $\varpi=\varpi'$ i.e. $F_{G_0}^f=F_{G'_0}^f$ (dans $A$). Dans ce cas, le groupe des homs en question, comme $A$-module, est isom au \deleted{\unsure{$E'$}} sous-module
\[
H\subset\operatorname{Hom}(\Omega',\Omega)
\]
formé des $v_0:\Omega'\longrightarrow\Omega$ tels que \inserted{pour tt $0\leq\nu\leq f-2$} on ait
\[
v_0\pi^{(\sum_0^\nu r_i)-(\sum_0^\nu r'_i)}\in\operatorname{Hom}(\Omega',\Omega)
\]
i.e.
\[
\operatorname{val}(v_0)+\sum_0^\nu r_i-\sum_0^\nu r'_i\geq0.
\]
